\documentclass[11pt]{article}
\usepackage[margin=1in]{geometry}
\usepackage{amsmath,amssymb,amsthm,booktabs,array}
\usepackage[T1]{fontenc}
\usepackage{lmodern}
\usepackage[colorlinks=true,linkcolor=blue,citecolor=blue,urlcolor=blue]{hyperref}
\newtheorem{theorem}{Theorem}
\newtheorem{lemma}{Lemma}
\newcommand{\R}{\mathbb R}
\newcommand{\C}{\mathbb C}
\newcommand{\rhoo}{\varrho}
\title{A rational counterexample to both Schur stability assertions\\
\large An exact specialization and Lean certificate for MathDB 370245}
\author{Myeongjin Kang}
\date{25 October 2026}
\setlength{\emergencystretch}{2em}
\begin{document}
\maketitle
\begin{abstract}
We verify a two-dimensional rational counterexample to Conjecture 1.1 of
Banerjee and Chaudhury, arXiv:2509.06715v1. A primitive row-stochastic matrix
$W$ and a positive-definite matrix $B$ meet every stated hypothesis, but both
$W(I-tB)$ and $I-W+(I+tB)^{-1}(2W-I)$ have spectral radius exactly one at
the strictly admissible step $t=1$. The example is the $s=1$ specialization
of Shivam Patel's previously posted family. No mathematical novelty or
priority is claimed. A single Lean module verifies the hypotheses, actual
inverse, complete complex spectra, and failure of both stability assertions
using mathlib's spectral radius.
\end{abstract}

\section{The conjecture and its precise scope}
For a real square matrix $A$, let $\sigma(A)$ be the spectrum of its
complexification and define
\[
 \rhoo(A)=\max\{|\lambda|:\lambda\in\sigma(A)\}.
\]
The matrix is Schur stable if $\rhoo(A)<1$. A matrix $W$ is row-stochastic
if its entries are nonnegative and $We=e$, where $e$ is the all-ones column
vector. It is primitive if $W^k$ has strictly positive entries for some
positive integer $k$. We use the normalized left Perron vector $\pi$:
$\pi>0$, $\pi^\top e=1$, and
$\pi^\top W=\rhoo(W)\pi^\top$.

Conjecture 1.1 of \cite{BC} asserts that, if $W$ is primitive and stochastic
and $B$ is positive semidefinite with
\begin{equation}\label{eq:hyp}
 Be\leq\rhoo(B)e,\qquad \pi^\top Be>0,
\end{equation}
then both operators
\begin{equation}\label{eq:operators}
 P(t)=W(I-tB),\qquad R(t)=I-W+(I+tB)^{-1}(2W-I)
\end{equation}
are Schur stable for every $0<t<2/\rhoo(B)$.
All vector inequalities here are componentwise.

This conjecture does not impose entrywise nonnegativity on $B$.
The paragraph immediately preceding it expressly discusses dropping that
condition; Example 1.2 gives a matrix with negative off-diagonal entries
and says that it satisfies the conjecture's assumptions. This distinction
matters because the abstract and some earlier results of \cite{BC} discuss
nonnegative matrices. Our target is the displayed Conjecture 1.1 with its
own hypotheses.

\clearpage
\section{The counterexample}
\begin{theorem}\label{thm:main}
Set
\begin{equation}\label{eq:data}
 W=\begin{pmatrix}0&1\\[1pt]1/4&3/4\end{pmatrix},\qquad
 B=\frac1{11}\begin{pmatrix}9&-4\\-4&3\end{pmatrix},\qquad
 t=1,\qquad \pi=\begin{pmatrix}1/5\\4/5\end{pmatrix}.
\end{equation}
Then $W$ is primitive and row-stochastic, $\pi$ is its normalized positive
left Perron vector, and $B$ is positive definite. Moreover,
\[
 \rhoo(B)=1,\qquad Be\leq\rhoo(B)e,\qquad
 \pi^\top Be=\frac1{55}>0,\qquad 0<t<\frac2{\rhoo(B)}.
\]
Nevertheless,
\[
 \sigma(P(1))=\{0,1\},\qquad \sigma(R(1))=\{1/2,1\}.
\]
Thus $\rhoo(P(1))=\rhoo(R(1))=1$, and neither operator is Schur stable.
Both assertions of Conjecture 1.1 are false.
\end{theorem}

\begin{proof}
The entries of $W$ are nonnegative and its row sums are one. Also,
\[
 W^2=\begin{pmatrix}1/4&3/4\\3/16&13/16\end{pmatrix}>0,
 \qquad
 \det(zI-W)=(z-1)(z+1/4).
\]
Consequently $W$ is primitive and $\rhoo(W)=1$. Direct multiplication gives
$\pi^\top W=\pi^\top$. Both coordinates of $\pi$ are positive and sum to
one, so this is the required normalized left Perron vector.

The matrix $B$ is symmetric. For $x=(x_1,x_2)^\top\in\R^2$,
\begin{equation}\label{eq:quadratic}
 11x^\top Bx=8(x_1-x_2/2)^2+x_1^2+x_2^2.
\end{equation}
For $x\ne0$ the last two terms have strictly positive sum, proving positive
definiteness and hence positive semidefiniteness. Its characteristic
polynomial factors as
\[
 \det(zI-B)=(z-1)(z-1/11).
\]
Equivalently, $(2,-1)^\top$ and $(1,2)^\top$ are eigenvectors with eigenvalues
$1$ and $1/11$, respectively. It follows that $\rhoo(B)=1$.
Furthermore,
\[
 Be=\frac1{11}\begin{pmatrix}5\\-1\end{pmatrix}\leq e,
 \qquad
 \pi^\top Be=\frac{5-4}{55}=\frac1{55}>0.
\]
Thus every condition in \eqref{eq:hyp} holds, and $0<1<2/\rhoo(B)=2$.

For the first operator, exact multiplication gives
\begin{equation}\label{eq:P}
 I-B=\frac1{11}\begin{pmatrix}2&4\\4&8\end{pmatrix},\qquad
 P(1)=\begin{pmatrix}4/11&8/11\\7/22&7/11\end{pmatrix}.
\end{equation}
Its trace is one and its determinant is zero, so
$\det(zI-P(1))=z(z-1)$. More concretely,
\[
 P(1)\begin{pmatrix}8\\7\end{pmatrix}
 =\begin{pmatrix}8\\7\end{pmatrix}.
\]
This nonzero eigenvector already obstructs strict Schur stability.

For the second operator, put
\[
 J=\begin{pmatrix}7/12&1/6\\1/6&5/6\end{pmatrix}.
\]
Direct multiplication verifies both $J(I+B)=I$ and $(I+B)J=I$;
therefore $J=(I+B)^{-1}$, the actual inverse appearing in
\eqref{eq:operators}. For additional transparency,
\[
 I-W=\begin{pmatrix}1&-1\\-1/4&1/4\end{pmatrix},\qquad
 J(2W-I)=\begin{pmatrix}-1/2&5/4\\1/4&3/4\end{pmatrix}.
\]
Adding them yields
\begin{equation}\label{eq:R}
 R(1)=\begin{pmatrix}1/2&1/4\\0&1\end{pmatrix}.
\end{equation}
This triangular matrix has characteristic polynomial $(z-1/2)(z-1)$;
indeed $R(1)(1,2)^\top=(1,2)^\top$.
The displayed factorizations hold for every complex $z$ and identify the
complete complex spectra, not just real roots. The two spectral radii are
therefore exactly one. This proves all assertions.
\end{proof}

The step $t=1$ lies strictly inside the conjectured open interval $(0,2)$.
The counterexample uses an eigenvalue on the unit circle, which suffices
because the conjecture requires a strict inequality. No eigenvalue
strictly larger than one, numerical root approximation, or limiting
argument is needed.

\section{Attribution and limits of the contribution}
Patel's reply \cite{Patel} already presents the family
\[
 W_s=\begin{pmatrix}0&1\\v&1-v\end{pmatrix},\qquad
 v=\frac1{5-s},\qquad 0<s<2,
\]
with the same $B$ as in \eqref{eq:data}. The posted calculation states that
both operators have eigenvalue one at $t=s$ and a real eigenvalue exceeding
one for $s<t<2$. Substituting $s=1$ gives exactly our $W$ and $t$; normalizing
$(v,1)^\top$ gives exactly our $\pi$.

Our contribution is an independent exact verification and a Lean
formalization of this specialization, supplied for the MathDB 370245
verification campaign. We do not claim discovery of the counterexample,
a new mathematical theorem, the first formalization, or a verification of
every assertion about the full posted family. The broader family is not
needed here. We do not challenge the separate local or special-case
results in \cite{BC}; failure at one admissible step does not contradict a
per-instance sufficiently-small-step theorem.

\section{Formal certificate and reproduction}
The accompanying module \nolinkurl{formal/SchurStability.lean} uses Lean 4.34.1
and pinned mathlib revision
\begin{center}\small\nolinkurl{d13f23b723b8a846827a245b89c10fc7d3f11612}\end{center}
It defines $W,B,t,\pi$ as the literal real data in \eqref{eq:data}, and defines
$P$ and $R$ by \eqref{eq:operators}. In particular, $R$ uses matrix inversion;
it is not defined by the simplified answer \eqref{eq:R}.

For a real matrix $A$, \texttt{complexify A} maps its entries through
$\R\hookrightarrow\C$. The definition \texttt{rho A} is mathlib's
\nolinkurl{spectralRadius} over $\C$ of this complex matrix. Its codomain is
the extended nonnegative reals. The definition \texttt{realRho A} takes
its real value. The four exact radius-one proofs in this example make
that conversion harmless; no infinite value is used as a real radius.

The spectrum calculation is connected to mathlib by the equivalence
\[
 z\in\sigma(A)\quad\Longleftrightarrow\quad \det(zI-A)=0
\]
for complex matrices. The library definition characterizes the left side
by noninvertibility; invertibility of a matrix is equivalent to
invertibility of its determinant. Over $\C$, the latter fails exactly when
the determinant is zero. The certificate proves this bridge, computes the
four factorizations, and uses
\[
 \operatorname{spectralRadius}_{\C}(A)
 =\sup_{z\in\sigma(A)}\|z\|
\]
to obtain the exact radii. It also formally identifies matrix spectra with
finite-dimensional linear-map eigenvalues, and verifies nonzero complex
eigenvectors of $P(1)$ and $R(1)$ with eigenvalue one.

\begin{center}
\small
\begin{tabular}{@{}>{\raggedright\arraybackslash}p{0.37\textwidth}>{\raggedright\arraybackslash}p{0.57\textwidth}@{}}
\toprule
Mathematical content & Lean declarations (namespace \nolinkurl{SchurStability})\\
\midrule
Stochasticity and primitivity & \nolinkurl{W_stochastic}, \nolinkurl{W_square}, \nolinkurl{W_primitive}\\
Perron vector & \nolinkurl{pi_positive}, \nolinkurl{pi_normalized}, \nolinkurl{pi_is_left_Perron}\\
Positive definiteness & \nolinkurl{B_quadratic}, \nolinkurl{B_posDef}, \nolinkurl{B_posSemidef}\\
Side conditions and strict step & \nolinkurl{B_source_bound}, \nolinkurl{pi_B_e_positive}, \nolinkurl{strict_step}\\
Actual inverse and operators & \nolinkurl{inverse_left}, \nolinkurl{inverse_right}, \nolinkurl{actual_inverse}, \nolinkurl{P_exact}, \nolinkurl{R_exact}\\
Complete complex spectra & \nolinkurl{B_spectrum}, \nolinkurl{W_spectrum}, \nolinkurl{P_spectrum}, \nolinkurl{R_spectrum}\\
Exact spectral radii & \nolinkurl{rho_B}, \nolinkurl{rho_W}, \nolinkurl{rho_P}, \nolinkurl{rho_R}\\
Eigenvector/eigenvalue bridge & \nolinkurl{eigenvalue_iff_spectrum}, \nolinkurl{P_complex_eigenvector}, \nolinkurl{R_complex_eigenvector}\\
All hypotheses and both failures & \nolinkurl{schur_counterexample}\\
Negated universal assertions & \nolinkurl{both_conjectures_false}\\
\bottomrule
\end{tabular}
\end{center}

The last declaration refutes both universal statements already in dimension
two, which is enough to refute their unrestricted versions. The source
hypotheses include normalization of the Perron vector, a harmless
strengthening of the source's choice of scaling.

From an extracted release directory, run
\begin{verbatim}
python3 verify_outcome.py .
\end{verbatim}
The campaign verifier checks pinned dependencies, compiles a byte-identical
isolated source copy with warnings treated as errors, and prints the type
and foundational dependencies of every declaration. The release includes the verifier and dependency pins. The completed
audit checks 77 declarations and reports only \nolinkurl{propext},
\nolinkurl{Classical.choice}, and \nolinkurl{Quot.sound}. An independent review accompanies the release; it is separate from kernel checking.

\begin{thebibliography}{9}
\bibitem{BC}
D. Banerjee and K. N. Chaudhury,
\emph{On the Schur Stability of Some Image Reconstruction Operators},
arXiv:2509.06715v1, 8 September 2025.
\url{https://arxiv.org/html/2509.06715v1}.
Conjecture 1.1, its preceding paragraph, and Example 1.2.

\bibitem{Patel}
S. Patel (username \nolinkurl{shivampatel}),
\emph{Counterexample to both stability assertions, with arbitrarily small
instability thresholds}, MathDB 370245, reply 1, 20 August 2026.
Frozen reply identifier
\nolinkurl{5a30dc24-61b4-49cc-a6de-ee5f67a71bf3},
recorded in \nolinkurl{metadata.json}.
\end{thebibliography}
\end{document}
